\documentclass[11pt]{article}
% Companion note to article.tex, English version; the French version is notes-fr.tex. Built with XeLaTeX:
% cd docs/bridges && xelatex notes.tex && xelatex notes.tex
\usepackage{polyglossia}\setmainlanguage{english}\setotherlanguage{french}
\newcommand{\thmname}{Theorem}\newcommand{\corname}{Corollary}\newcommand{\propname}{Proposition}
\newcommand{\lemname}{Lemma}\newcommand{\bridgename}{Bridge}\newcommand{\defname}{Definition}
\newcommand{\exname}{Example}\newcommand{\exoname}{Exercise}\newcommand{\remname}{Remark}
% Shared preamble of docs/bridges/article.tex (English) and article-fr.tex (French).
% Built with XeLaTeX: cd docs/bridges && xelatex article.tex && xelatex article.tex
\usepackage{fontspec}
\IfFontExistsTF{Charter}{\setmainfont{Charter}}{\IfFontExistsTF{XCharter}{\setmainfont{XCharter}}{}}
\setmonofont{DejaVu Sans Mono}[Scale=0.8]
\usepackage{microtype}
\usepackage[a4paper,margin=27mm]{geometry}
\usepackage{setspace}\setstretch{1.1}
\usepackage{amsmath,amssymb,amsthm}
\usepackage{mathtools}
\usepackage[most]{tcolorbox}
\usepackage{xcolor}
\usepackage{tikz}
\usetikzlibrary{arrows.meta,positioning,calc,cd}
\usepackage{booktabs,array,tabularx}
\usepackage{fancyvrb}
\usepackage[hidelinks]{hyperref}
\usepackage{xurl}
\usepackage{enumitem}
\setlist{nosep,leftmargin=*}
\setlength{\emergencystretch}{2em}
\setlength{\parskip}{0.35em}
\numberwithin{equation}{section}

\newcommand{\Set}{\mathbf{Set}}
\newcommand{\PSh}{\mathrm{PSh}}
\newcommand{\Sh}{\mathrm{Sh}}
\newcommand{\Kar}{\operatorname{Kar}}
\newcommand{\Hom}{\operatorname{Hom}}
\newcommand{\End}{\operatorname{End}}
\newcommand{\Aut}{\operatorname{Aut}}
\newcommand{\id}{\mathrm{id}}
\newcommand{\B}{\mathbf{B}}
\newcommand{\y}{\mathbf{y}}
\newcommand{\op}{\mathrm{op}}
\newcommand{\res}[2]{#1|_{#2}}
\newcommand{\q}[1]{«\,#1\,»}
\newcommand{\lean}[1]{{\small\nolinkurl{#1}}}
\newcommand{\lp}{\href{https://github.com/LeanPool/lean-pool/tree/main/LeanPool/ToposTheoreticBridges}{\nolinkurl{ToposTheoreticBridges}}}

% Lean code blocks.
\DefineVerbatimEnvironment{leancode}{Verbatim}{fontsize=\small,frame=leftline,framerule=0.6pt,rulecolor=\color{blue!40},xleftmargin=6pt}

\theoremstyle{plain}
\newtheorem{theorem}{\thmname}[section]
\newtheorem{corollary}[theorem]{\corname}
\newtheorem{proposition}[theorem]{\propname}
\newtheorem{lemma}[theorem]{\lemname}
\newtheorem{bridge}{\bridgename}
\renewcommand{\thebridge}{\Alph{bridge}}
\theoremstyle{definition}
\newtheorem{definition}[theorem]{\defname}
\newtheorem{example}[theorem]{\exname}
\newtheorem{exercise}{\exoname}[section]
\theoremstyle{remark}
\newtheorem*{remark}{\remname}
\tcolorboxenvironment{theorem}{enhanced, breakable, colback=green!2, colframe=green!35!black, boxrule=0.6pt, arc=0pt, left=7pt, right=7pt, top=4pt, bottom=4pt, before skip=12pt, after skip=8pt}
\tcolorboxenvironment{bridge}{enhanced, breakable, colback=blue!3, colframe=blue!35!black, boxrule=0.5pt, arc=0pt, left=7pt, right=7pt, top=4pt, bottom=4pt, before skip=12pt, after skip=8pt}
\tcolorboxenvironment{definition}{enhanced, breakable, colback=black!3, colframe=black!25, boxrule=0.4pt, arc=0pt, left=7pt, right=7pt, top=4pt, bottom=4pt, before skip=10pt, after skip=6pt}
\newcommand{\entry}[1]{\par\smallskip\noindent\textbf{#1}\enspace}

\newcommand{\inwords}{\entry{In words.}}
\newcommand{\tracename}{Pipeline trace}
\newcommand{\gen}{\entry{How the generator finds it.}}
\newcommand{\direct}{\entry{Proof, without topoi.}}
\newcommand{\status}{\entry{Is it new?}}
\newcommand{\leanentry}{\entry{Lean.}}

\title{Bridges in Grothendieck's notes}
\author{Michel Hua\\ \small Commutator, Paris\\ \small\texttt{michel@commutator.io} · ORCID 0009-0002-8304-0683}
\date{Draft, 23 November 2026}

\begin{document}
\maketitle
\begin{abstract}
O.~Caramello's \q{bridges} read one property of a topos on two presentations of it and turn the two readings into a theorem relating the presentations. A companion article \cite{bridges-generator} runs such bridges by hand as a theorem generator. This note applies the same generator to Grothendieck's manuscript notes (Archives Grothendieck, Université de Montpellier). Three arches (two-valued, connected, presheaf topos of a preorder), proved in Lean~4 in this edition's package \texttt{lean/bridges}, are read on three decks that recur in his notes: the simplex category, groups and preorders. We discuss two precedents in his notes: an exercise of folder 161-3, p.~26, which the generator re-derives from the Morita arch and whose preorder form is proved in Lean, with the same identification in folder 37, p.~48; and the correspondence between subtopoi and quotients of a theory in folder 161-2, pp.~17 and~44. A last deck from folder 161-2, the topos of pseudo-fields (von Neumann regular rings), gets a section of its own: we check one of his three presentations of it, read the arches on it and answer a question he left in the margin. An appendix lists the Lean files. A French version accompanies this text.
\end{abstract}

\tableofcontents
\section{Introduction}\label{sec:nintro}

A topos can be presented in many ways that look unrelated: as the actions of a monoid, as the representations of a quiver, as the sheaves on a space, as the models of a logical theory. O.~Caramello's \emph{bridges} \cite{caramello-unification,caramello-tst} exploit this. A property of the topos that is invariant under equivalence is characterised on each of two presentations; these characterisations are the \emph{arches} of the bridge, and the equivalence between the two presentations is its \emph{deck}. Put together, the two arches give a theorem relating the two presentations, in which the topos no longer appears.

In a talk at Mines Paris \cite{caramello-mines}, Caramello proposed that a machine combine arches and decks systematically. The companion article \cite{bridges-generator} runs a small version of this generator by hand and describes it as a pipeline: a registry of arches, proved in Lean~4 as lemmas of the form \texttt{(property of PSh C) ↔ (condition on C)}; a registry of decks, small categories $C$ whose presheaves are a familiar structure; and stages that pair every deck with every arch, compute the condition on $C$, rewrite it in the language of the field, and check it in Lean, by a direct proof and against the literature. This note applies the generator to Grothendieck's manuscript notes kept at the University of Montpellier, which already contain decks of this kind and invariants read on two presentations of one topos.

\subsection*{Plan}

Section~\ref{sec:recall} recalls the notation and the five arches of the companion article. Section~\ref{sec:morearches} gives three more arches (two-valued, connected, spatial), proved in Lean, and Section~\ref{sec:grdecks} reads them on three decks that recur in Grothendieck's notes. Section~\ref{sec:precedents} collects two precedents in his notes: the exercise of folder 161-3, p.~26, which the generator re-derives and whose preorder form we prove in Lean, with the same identification in folder 37, p.~48; and the subtopoi of a theory in folder 161-2. Section~\ref{sec:pseudofields} adds a deck taken from folder 161-2, the topos of pseudo-fields, which Grothendieck presents in three ways; we check one of the presentations, read the arches on the deck, and answer a question he left in the margin. Section~\ref{sec:stacks} relates the classifying topoi to Grothendieck's stacks and proposes finer arches, with an example from folder 146. Section~\ref{sec:nconclusion} concludes, and Appendix~\ref{sec:nlean} lists the Lean files and declarations.

\section{Notation and the five arches}\label{sec:recall}

We use the notation of the companion article \cite{bridges-generator}, recalled here briefly. A small category $C$ is a \emph{shape}. A \emph{presheaf} on $C$ is a functor $C^{\op}\to\Set$, in Lean \texttt{Cᵒᵖ ⥤ Type u}; presheaves on $C$ form the topos $\PSh(C)$, which Grothendieck writes $\widehat C$. A \emph{deck} is a shape whose presheaves are a familiar structure. The negation $\neg A$ of a sub-presheaf $A\subseteq F$ is the largest sub-presheaf of $F$ disjoint from $A$. $\PSh(C)$ is \emph{Boolean} if $A\cup\neg A = F$ for every $A\subseteq F$ (excluded middle), and \emph{De Morgan} if $\neg A\cup\neg\neg A = F$ for every $A\subseteq F$ (weak excluded middle). $C$ satisfies the \emph{Ore condition} if every cospan $f : Y\to X$, $g : Z\to X$ completes to a commutative square $f\circ h = g\circ k$. The \emph{Karoubi envelope} $\Kar(C)$ adds an object for each idempotent arrow $e = e\circ e$, and contains $C$ as a full subcategory. $C$ is \emph{thin} if there is at most one arrow between any two objects.

The companion article runs its generator on five arches, proved in Lean~4 in the library \lp{} of lean-pool:
\begin{itemize}
\item \emph{Boolean:} $\PSh(C)$ is Boolean iff $C$ is a groupoid;
\item \emph{De Morgan:} $\PSh(C)$ is De Morgan iff $C$ satisfies the Ore condition;
\item \emph{Morita:} $\PSh(C)\simeq\PSh(D)$ iff $\Kar(C)$ and $\Kar(D)$ are equivalent;
\item \emph{essential points:} the functors $\PSh(C)\to\Set$ that preserve all limits and all colimits form a category equivalent to $\Kar(C)^{\op}$;
\item \emph{subtopoi:} the subtopoi of $\PSh(C)$ correspond, order-reversingly, to the Grothendieck topologies on $C$.
\end{itemize}

\section{Three more arches}\label{sec:morearches}

We add three arches to the registry of the companion article \cite{bridges-generator}. So far we have applied them only to the decks taken from Grothendieck's notes (\S\ref{sec:grdecks}) and to his pseudo-fields (\S\ref{sec:pseudofields}); none of the outputs of the companion article comes from them. We number them F to H, after the five arches A to E of the companion article.
\setcounter{bridge}{5}

\begin{bridge}[Two-valued]\label{br:twoval}
The subterminal objects of $\PSh(C)$ (its \q{opens}) correspond to the sieves of $C$ \cite[IV~8.4.4]{sga4}. Hence $\PSh(C)$ has no subterminal objects besides $0$ and $1$ iff for any two objects $a$, $b$ of $C$ there is an arrow $a\to b$. For $C$ empty the topos is degenerate and the condition holds vacuously, so the statement is about nonempty $C$. Lean: \lean{forall_subfunctor_terminal_eq_bot_or_eq_top_iff}.
\end{bridge}

\begin{bridge}[Connected]\label{br:conn}
$\PSh(C)$ is connected and nonempty in the sense of \cite[IV~8.7~e)]{sga4} (every partition of its final object has exactly one nonempty part) iff $C$ is connected, a fact recalled as known in \cite[Ex.~3.6]{hora-kamio}. In terms of subterminal objects: $C$ is nonempty and the only complemented subterminal objects are $0$ and $1$. Lean: \lean{isConnected_iff_forall_isCompl}.
\end{bridge}

\begin{bridge}[Spatial]\label{br:spatial}
$\PSh(C)$ is equivalent to the presheaf topos of a preorder, and so to the topos of sheaves on its Alexandrov space, iff $C$ is thin. Grothendieck states it as an exercise \cite[p.~26]{grothendieck-161-3}; see \S\ref{sec:precedents}. Lean: \lean{exists_preorder_moritaEquivalentCat_iff_isThin}, which states the preorder form; that the topos is then the topos of a space is Grothendieck's equivalence, from the same exercise, and is not formalised here.
\end{bridge}

On the two main decks of the companion article, right $M$-sets and representations of a quiver, they read as follows.

\begin{center}\small
\begin{tabularx}{\linewidth}{>{\raggedright\arraybackslash}p{2.6cm}>{\raggedright\arraybackslash}p{3.2cm}>{\raggedright\arraybackslash}X>{\raggedright\arraybackslash}X}
\toprule
invariant & reading on $C$ & on right $M$-sets & on representations of a quiver\\
\midrule
two-valued & an arrow between any two objects & always & strongly connected\\
connected & $C$ connected & always & connected\\
preorder (spatial) & $C$ thin & $M$ trivial & at most one path between two vertices\\
\bottomrule
\end{tabularx}
\end{center}

\section{The decks of Grothendieck's notes}\label{sec:grdecks}

Three decks recur in Grothendieck's working notes. Here is what the three arches of \S\ref{sec:morearches} say about them (Lean: \lean{GrothendieckDecks.lean}).
\begin{itemize}
\item The simplex category $\Delta$, the first test category of \emph{Pursuing Stacks} (folders 108 and 118 of the fonds): simplicial sets form a two-valued and connected topos, since there is a constant map between any two objects of $\Delta$; this topos is not the presheaf topos of a preorder, since the two vertices of $[1]$ are two maps $[0]\to[1]$ (\lean{simplexCategory_twoValued}, \lean{simplexCategory_isConnected}, \lean{simplexCategory_not_spatial}).
\item A group $G$ as a one-object groupoid, whose presheaf topos is $\B G$ (folder 161-6, p.~22, where Grothendieck reduces a question to this case): the topos is two-valued (as \cite[IV~8.4.3]{sga4} notes for any monoid) and connected, and it is the presheaf topos of a preorder iff $G$ is trivial (\lean{singleObj_twoValued}, \lean{singleObj_isConnected}, \lean{singleObj_spatial_iff}).
\item A preordered set $I$, whose presheaf topos $\widehat I$ is $\mathrm{Top}(I)$ (folder 37, p.~48; folder 161-3, p.~26): the topos is two-valued iff any two elements are comparable both ways, that is, iff its posetal quotient has at most one point (\lean{preorder_twoValued_iff}).
\end{itemize}
These are immediate corollaries. We record them because they are the first entries of the registry taken from the fonds.

\section{Precedents in Grothendieck's notes}\label{sec:precedents}

This section collects two precedents in Grothendieck's notes. The decks taken from his notes are already in the registry (\S\ref{sec:grdecks}).

\subsection{Preorders and spatial topoi: folder 161-3, p.~26, and folder 37, p.~48}

In an undated manuscript note on topoi \cite{grothendieck-161-3}, p.~26, Grothendieck lists equivalent conditions on a topos $E$: to be equivalent both to the topos of a space and to a presheaf topos; to be the topos of an \q{essential} space (a space all of whose points are essential); to be equivalent to $\widehat I$ for an ordered set $I$; to be generated by the subobjects of the final object and to have enough essential points. He asks for a proof that $\widehat C$ satisfies these conditions iff $C$ is preordered, and that the topos of a space $X$ does iff the essential points of $X$ are very dense. He also notes 2-equivalences between essential spatial topoi, preorders and essential spaces. In our terms these are bridges on the deck of preorders, with the arches of subterminal objects and of essential points. The generator re-derives the first exercise from the Morita arch: if $\widehat C\simeq\widehat P$ for a preorder $P$, then $\Kar(C)\simeq\Kar(P)\simeq P$, and $C$, a full subcategory of its Karoubi envelope, is thin; conversely a thin category is equivalent to its thin skeleton, a partial order (Lean: \lean{exists_preorder_moritaEquivalentCat_iff_isThin}, below). The second exercise, on spaces, is not formalised here. The same identification appears in his notes for the re-edition of SGA \cite[p.~48]{grothendieck-37}: for an ordered set $I$ the opens of $\widehat I$ are the sieves of $I$, hence $\widehat I\simeq\mathrm{Top}(I)$. This step uses that $\widehat I$ is generated by its subterminal objects, which holds because for a preorder the representable presheaves are subterminal.

The Lean statement of the first exercise, in \lean{PreorderPresheaves.lean}, is the following. Its proof follows the argument above: \lean{karoubi_of_moritaEquivalentCat} reads the Morita arch, \lean{karoubiEquivPreorder} identifies a preorder with its Karoubi envelope, and the converse goes through the thin skeleton.
\begin{leancode}
variable (C : Type u) [SmallCategory C]

theorem exists_preorder_moritaEquivalentCat_iff_isThin :
    (∃ (P : Type u) (_ : Preorder P), MoritaEquivalentCat C P) ↔
      Quiver.IsThin C
\end{leancode}
The companion article lists this statement in its calibration log, as not in print as such \cite{bridges-generator}.

\subsection{Subtopoi and quotients of a theory: folder 161-2, pp.~17 and 44}

In his notes on universal algebra \cite[p.~17]{grothendieck-161-2}, under the heading \q{sous-théorie d'une $\Delta$-théorie (représentable par un topos) et sous-topos}, Grothendieck takes a topos $\mathcal T$ (his $R$) and a set $\Sigma$ of arrows of $\mathcal T$: the morphisms that invert the arrows of $\Sigma$ factor through a subtopos $\mathcal T'$ of $\mathcal T$, and every subtopos of $\mathcal T$ arises in this way. When $\mathcal T$ is the classifying topos of a theory, this is in substance the correspondence between quotients of the theory and subtopoi of its classifying topos on which Caramello builds her bridges \cite{caramello-tst}. On p.~44 of the same folder he reads structures as open subtopoi of a classifying topos.

\section{A deck from Grothendieck's notes: pseudo-fields}\label{sec:pseudofields}

The precedents of \S\ref{sec:precedents} show Grothendieck reading invariants of a topos on two presentations. His notes on universal algebra also contain a deck that the registry did not have: the topos of \q{pseudo-fields}, for which he gives three presentations.

\subsection{The deck}

In folder 161-2 \cite{grothendieck-161-2}, pp.~38--43, Grothendieck lists properties of rings and, for each, a subtopos of the topos $X_{\mathrm{ann}}\simeq\widehat R$ classifying rings, where $R$ is the category of affine schemes of finite type over $\mathbb Z$. His numbering changes between two pages: on p.~41 item~12 is the fields, while on p.~42 he renumbers, so that item~12 becomes the compact rings and item~13 the compact reduced rings. Below, \q{item~12} and \q{item~13} refer to the numbering of p.~42, and the fields are \q{the fields of p.~41}. Item~13 is the \emph{compact reduced ring}, which he names \emph{pseudo-field}: for every $f$ there is a unique idempotent $e$ with $ef$ invertible in $eA$ and $(1-e)f=0$, so that $A\simeq A_f\times A/fA$. These are the commutative von Neumann regular rings. He gives three presentations of the subtopos they define:
\[
X_{\mathrm{ps.\,corps}}\;\simeq\;\widetilde R_{\mathrm{ps.\,corps}}\;\simeq\;\widehat{R_{\mathrm{cons}}}.
\]
The second is the topos of sheaves on $R$ for the topology in which a family covers when one of its members has a section fibre by fibre; the third is the topos of presheaves on the category $R_{\mathrm{cons}}$ obtained from $R$ by inverting the arrows that are bijective with trivial residue field extensions (pp.~41--42).

\subsection{Checking the presentation by presheaves}

The theory of pseudo-fields can be axiomatised by the axioms of commutative rings and the sequent
\[
\forall a\ \exists!\,b\ \ (a^2b = a\ \wedge\ ab^2 = b):
\]
in a commutative ring the quasi-inverse $b$ of $a$ is unique, and the idempotent of Grothendieck's definition is $e = ab$. So the theory is cartesian, hence of presheaf type: its classifying topos is the topos of $\Set$-valued functors on its finitely presentable models. Let $T$ be the left adjoint of the inclusion of von Neumann regular rings into rings. This inclusion preserves filtered colimits, so $T$ preserves finite presentability. Conversely, every finitely presentable pseudo-field is some $T(A)$ with $A$ a finitely presented ring: in a finite presentation, replace each quasi-inverse by a new variable subject to its two equations. By Olivier \cite{olivier}, $\operatorname{Spec}T(A)$ is $\operatorname{Spec}A$ with its constructible topology, and the local rings of $T(A)$ are the residue fields of $A$. By Pierce \cite{pierce}, a von Neumann regular ring is the ring of global sections of the sheaf of its local rings over its spectrum, a Boolean space. Hence for $u:A\to B$ of finite presentation, $T(u)$ is an isomorphism iff $\operatorname{Spec}u$ is bijective and every residue field extension is trivial: a continuous bijection of Boolean spaces is a homeomorphism, and a map of sheaves that is an isomorphism on stalks is an isomorphism. The arrows that $T$ inverts are therefore exactly the arrows of the class $W$ that Grothendieck inverts.

It remains to see that the finitely presentable pseudo-fields are obtained from the finitely presented rings by inverting $W$, and we argue as follows. Olivier's construction writes $T(B)$ as a filtered colimit of finitely presented $B$-algebras obtained by adjoining quasi-inverses, one element $a$ at a time:
\[
B' = B[y]/(a^2y - a,\ ay^2 - y).
\]
Each map $B\to B'$ of this tower is bijective on spectra with trivial residue field extensions: over a point where $a$ is nonzero in the residue field, the equations force $y = 1/a$, and where $a = 0$ they force $y = 0$. So these maps lie in $W$, and so do their composites. Since $A$ is finitely presented, a map $T(A)\to T(B)$, that is a map $A\to T(B)$, factors through a stage of the tower, and is therefore represented by a roof $A\to B'\leftarrow B$ with $B\to B'$ in $W$. Hence the category of finitely presentable pseudo-fields is the localisation of the category of finitely presented rings at $W$; its opposite is $R_{\mathrm{cons}}$, and the classifying topos is $\widehat{R_{\mathrm{cons}}}$, as Grothendieck states. We have not checked the presentation by the topology $\mathrm{ps.\,corps}$. Nothing in this section is formalised.

\subsection{The arches on this deck}

\begin{itemize}
\item \emph{De Morgan:} yes; the finitely presentable models of a cartesian theory have pushouts, so they amalgamate (the De Morgan arch, \S\ref{sec:recall}).
\item \emph{Two-valued:} no; the sentence $1=0$ defines a subterminal object that is neither $0$ nor $1$.
\item \emph{Connected:} yes; $T(\mathbb Z)$ is an initial finitely presentable model, so the category of these models is connected.
\item \emph{Boolean:} no; the map $T(\mathbb Z)\to T(\mathbb F_2)$ is not invertible.
\item \emph{Essential points:} the finitely presentable pseudo-fields $T(A)$, up to isomorphism.
\item \emph{Subtopoi:} the quotient theories of the theory of pseudo-fields, by the duality between quotient theories and subtopoi \cite{caramello-tst}. Several items of Grothendieck's list are such quotients: the fields of p.~41, which he identifies with the integral pseudo-fields and with the local ones (he marks the last equality \q{à prouver}, and it is true, since a local von Neumann regular ring is a field); separably and algebraically closed fields (items~14--15); rings killed by $n$ and rings in which $n$ is invertible (item~16, p.~43), which he marks as open subtopoi.
\end{itemize}

\subsection{An answer to a question in the margin}

For item~12, the \emph{compact} rings (every $f$ invertible on an idempotent $e$ and nilpotent on $1-e$, p.~42), Grothendieck notes in the margin that the topology is not coherent and asks whether the topos has enough points. On p.~41 he also warns that the subtopoi of his list are determined by their points, but that this is not clear for their intersections, so that checking an inclusion between intersections on set-theoretic structures may not suffice. We settle both doubts through axioms. Each item of his list is axiomatised over the countable signature of rings by finitely many geometric sequents. For the compact rings one of these sequents has a countable disjunction, since $(1-e)f$ is nilpotent of some order $n$ with no bound on $n$; this is why the topology is not coherent. The classifying topos of a countable geometric theory is the topos of sheaves on its syntactic site, a countable category with finite limits whose topology is generated by countably many sieves. A theorem of Makkai and Reyes \cite[Thm~6.2.4]{makkai-reyes} (see also \cite[Cor.~3.3.10]{diliberti-rogers}) says that a topos of sheaves on a countable category with pullbacks, for a topology generated by a countable family of sieves, has enough points. So the topos of compact rings has enough points; the topos of pseudo-fields has enough points trivially, being a presheaf topos. Once a subtopos $E'$ has enough points, an inclusion $E'\subseteq E''$ can be checked on points. An intersection of countably many subtopoi of his list classifies the theory obtained by putting their axioms together, which is still countable, so it has enough points and his checks on points are legitimate for it. The argument does not cover intersections of uncountably many subtopoi.

\section{Stacks and finer arches}\label{sec:stacks}

The arches used so far answer yes or no: Boolean or not, connected or not. Stacks, Grothendieck's \emph{champs}, suggest arches that compute.

\subsection{A classifying topos represents a stack}

The classifying topos $\Set[\mathbb T]$ of a geometric theory $\mathbb T$ represents, on the 2-category of Grothendieck topoi, the pseudo-functor that sends a topos $\mathcal E$ to the category of models of $\mathbb T$ in $\mathcal E$ \cite{caramello-tst}: to classify structures is to represent a stack of models. The simplest instance is the deck $\B G$ of a group, read in \S\ref{sec:grdecks}: its topos classifies the $G$-torsors, so it is the stack of torsors, and the gerbes of Giraud, classified by a non-abelian $H^2$ \cite{giraud}, continue the same line.

\subsection{Quotient stacks as decks with two presentations}

A quotient stack $[X/G]$ comes with two presentations: by the groupoid $X\times G\rightrightarrows X$, whose topos is that of $G$-equivariant sheaves on $X$, and as a moduli problem, a stack of families of some kind of object. Reading an invariant on both is a bridge in the sense of this note, and the invariants that suit stacks are numerical or algebraic rather than logical: the Picard group, cohomology, the points. Mumford's computation $\operatorname{Pic}(\mathcal M_{1,1})\cong\mathbb Z/12$ \cite{mumford} is the archetype: an invariant of a moduli problem obtained from an explicit presentation.

\subsection{An example from the fonds}

In folder 146 Grothendieck tabulates, for $n$ marked points on the projective line, the order of the tangent line at the last point, a line bundle on $M_{0,n}$ that descends to the quotient by $\mathfrak S_{n-1}$; he gives the orders for four points and only guesses for five and six. The edition's article on this table \cite{hua-tangent} reads the Picard group on a second presentation of the stack: $[M_{0,n}/\mathfrak S_{n-1}]$ is the stack of monic squarefree polynomials of degree $n-1$ modulo the affine group, and on that presentation $\operatorname{Pic}[M_{0,n}/\mathfrak S_{n-1}]$ comes out cyclic of order $(n-1)(n-2)$, generated by the tangent line, in characteristic zero and for every $n\ge3$. For even $n\ge6$ the group was known through the integral Chow ring computed by Edidin and Hu \cite{edidin-hu}. This is a bridge with a Picard arch, and it gives a statement that our searches did not find in print for odd $n$; the proof and its Lean formalisation are in \cite{hua-tangent}.

\subsection{Stacks in the homotopical sense}

\emph{À la poursuite des champs} \cite{maltsiniotis} is first of all about stacks in the homotopical sense, and its test categories are bridges of yet another kind: a category $A$ is a test category when presheaves on $A$ model homotopy types, and the conditions Grothendieck gives (asphericity conditions on $A$, the existence of a separating interval in $\widehat A$) read a homotopical property of $\widehat A$ on the category $A$. The deck $\Delta$ of \S\ref{sec:grdecks} is the first test category. Folders 108 (a list of questions), 118 (scratchwork on the homotopy theory of $(\mathrm{Cat})$) and 134-2 (the typescript of the first chapter of \emph{Pursuing Stacks}, with its tables of contents) work on these conditions; \S\ref{sec:test} takes up their open questions.

\subsection{A programme}

These examples suggest a family of finer arches for the generator: for a quotient stack $[X/G]$, or more generally a stack with an explicit groupoid presentation, the Picard group, the cohomology and the points read on the presentation and compared with the moduli description. Unlike the arches of \S\ref{sec:morearches}, they produce numbers and groups, so their outputs can be tested against tables, such as Grothendieck's in folder 146, and against the OEIS. We have not run the generator with them; this section only records the connection and one worked example.

\section{Conclusion}\label{sec:nconclusion}

On the decks of Grothendieck's notes, the three arches of \S\ref{sec:morearches} give immediate corollaries; we record them because they are the first entries of the registry taken from the fonds. The note of folder 161-3, p.~26, contains two exercises that are bridges on the deck of preorders: the generator re-derives the first from the Morita arch, and its preorder form is proved in Lean; the second, on spaces, is not formalised. The pseudo-fields are the one deck of the registry between non-trivially equivalent theories. We checked their presentation by presheaves by hand, not their presentation by the topology $\mathrm{ps.\,corps}$, and nothing about them is formalised; the argument on points covers intersections of countably many subtopoi of his list, not of uncountably many. The depth of a generated theorem is bounded by the depth of its arch and of its deck \cite[slide 8]{caramello-mines}, and deeper readings of these notes would need arches on genuine sites, which the registry does not contain.

\appendix
\section{Lean}\label{sec:nlean}

The three arches of \S\ref{sec:morearches} are proved in the Lake package \texttt{lean/bridges} of this edition's repository (\url{https://github.com/Commutator-IO/grothendieck-project}), which depends on lean-pool at a fixed revision; the corollaries on the decks of Grothendieck's notes are formalised there too. The five arches of \lp{}, in lean-pool, are proved in Lean~4 with Mathlib, with no axioms beyond \texttt{propext}, \texttt{Classical.choice} and \texttt{Quot.sound}. This statement describes the development; the files contain no \texttt{\#print axioms} command, and we do not record a separate check for them. For \texttt{lean/bridges}, the continuous integration of this edition's repository runs an axiom audit on every change: at the version described here it covered 437 declarations, before the file \lean{TestCategories.lean} of \S\ref{sec:test} was added, all within \texttt{propext}, \texttt{Classical.choice} and \texttt{Quot.sound}. The files of the package used here are the following; the others serve the companion article \cite{bridges-generator}.
\begin{itemize}
\item \lean{PreorderPresheaves.lean} (\lean{exists_preorder_moritaEquivalentCat_iff_isThin}, after Grothendieck's exercise in folder 161-3, p.~26); \lean{SubterminalArches.lean} (the two-valued and connected arches: \lean{forall_subfunctor_terminal_eq_bot_or_eq_top_iff}, \lean{isConnected_iff_forall_isCompl}); \lean{GrothendieckDecks.lean} (the decks of Grothendieck's notes: \lean{simplexCategory_twoValued}, \lean{simplexCategory_isConnected}, \lean{simplexCategory_not_spatial}, \lean{singleObj_twoValued}, \lean{singleObj_isConnected}, \lean{singleObj_spatial_iff}, \lean{preorder_twoValued_iff}).
\end{itemize}
Not formalised: the equivalence between the presheaf topos of a preorder and the topos of a space (Bridge~\ref{br:spatial}), the second exercise of folder 161-3, p.~26 (\S\ref{sec:precedents}), and the statements of \S\ref{sec:pseudofields}.

\begin{center}\scriptsize\emergencystretch=3em
\begin{tabularx}{\linewidth}{>{\raggedright\arraybackslash}p{2.9cm}>{\raggedright\arraybackslash}X>{\raggedright\arraybackslash}p{2.4cm}>{\raggedright\arraybackslash}p{2.2cm}}
\toprule
statement & Lean declaration & file & status\\
\midrule
small categories: presheaf topos of a preorder iff thin (Grothendieck, 161-3, p.~26) & \lean{exists_preorder_moritaEquivalentCat_iff_isThin} & \lean{PreorderPresheaves} & proved in Lean\\
two-valued and connected arches (Bridges~\ref{br:twoval} and~\ref{br:conn}) & \lean{forall_subfunctor_terminal_eq_bot_or_eq_top_iff}, \lean{isConnected_iff_forall_isCompl} & \lean{SubterminalArches} & proved in Lean\\
decks of Grothendieck's notes: $\Delta$, groups, preorders (\S\ref{sec:grdecks}) & \lean{simplexCategory_twoValued}, \lean{simplexCategory_isConnected}, \lean{simplexCategory_not_spatial}, \lean{singleObj_twoValued}, \lean{singleObj_isConnected}, \lean{singleObj_spatial_iff}, \lean{preorder_twoValued_iff} & \lean{GrothendieckDecks} & proved in Lean\\
test categories and the De Morgan arch (\S\ref{sec:test}) & \lean{rightOreCondition_iff_pullback_ne_bot}, \lean{not_isPreconnected_elements}, \lean{not_isConnected_elements_sieves}, \lean{not_isConnected_elements_yonedaProdSieves}, \lean{not_isConnected_elements_sieves_of_isDeMorganTopos}, \lean{not_isConnected_elements_sieves_simplexCategory_op}, \lean{not_isConnected_elements_sieves_singleObj}, \lean{not_isConnected_elements_sieves_preorder}, \lean{not_isConnected_elements_sieves_int}, \lean{not_isConnected_elements_sieves_multiplicative_nat}, \lean{not_isDeMorganTopos_of_isConnected_elements_sieves} & \lean{TestCategories} & proved in Lean (connectedness core)\\
\bottomrule
\end{tabularx}
\end{center}

\begin{thebibliography}{99}
\bibitem{caramello-mines} O.~Caramello, \emph{Topos-empowered automatic theorem generation}, talk, CRI, Mines Paris, 17 September 2024. \url{https://www.oliviacaramello.com/Talks/MinesTalk.pdf}
\bibitem{grothendieck-161-3} A.~Grothendieck, \emph{Topos : notes manuscrites} (undated), Archives Grothendieck, Université de Montpellier, cote 161-3, \url{https://grothendieck.umontpellier.fr/}.
\bibitem{sga4} M.~Artin, A.~Grothendieck, J.-L.~Verdier, \emph{Théorie des topos et cohomologie étale des schémas} (SGA~4), tome~1, Lecture Notes in Math.~269, Springer, 1972; Exposé IV, A.~Grothendieck and J.-L.~Verdier, \emph{Topos}; re-typeset edition \url{https://www.normalesup.org/~forgogozo/SGA4/}.
\bibitem{grothendieck-161-2} A.~Grothendieck, \emph{Algèbre universelle [ou catégories] : notes manuscrites} (undated), Archives Grothendieck, Université de Montpellier, cote 161-2, \url{https://grothendieck.umontpellier.fr/}.
\bibitem{grothendieck-37} A.~Grothendieck, \emph{Rééditions SGA, EGA}, notes and typescripts (1967--1971), Archives Grothendieck, Université de Montpellier, cote 37, \url{https://grothendieck.umontpellier.fr/}.
\bibitem{elephant} P.~T. Johnstone, \emph{Sketches of an Elephant: A Topos Theory Compendium}, vols.~1--2, Oxford Logic Guides 43--44, Oxford University Press, 2002.

\bibitem{caramello-flat} O.~Caramello, Extensions of flat functors and theories of presheaf type, arXiv:1404.4610.

\bibitem{caramello-johnstone} O.~Caramello and P.~T. Johnstone, De Morgan's law and the theory of fields, arXiv:0808.1972.

\bibitem{caramello-lattices} O.~Caramello, Lattices of theories, arXiv:0905.0299.

\bibitem{hutzler} M.~Hutzler, Syntactic presentations for glued toposes and for crystalline toposes, arXiv:2206.11244.

\bibitem{olivier} J.-P.~Olivier, \emph{Anneaux absolument plats universels et épimorphismes à buts réduits}, Séminaire Samuel, Algèbre commutative 2 (1967--1968), exposé~6 (Zbl 0162.05701).
\bibitem{pierce} R.~S.~Pierce, \emph{Modules over commutative regular rings}, Mem. Amer. Math. Soc.~70, 1967 (Zbl 0152.02601).
\bibitem{makkai-reyes} M.~Makkai, G.~E.~Reyes, \emph{First Order Categorical Logic. Model-theoretical methods in the theory of topoi and related categories}, Lecture Notes in Math.~611, Springer, 1977 (Zbl 0357.18002).
\bibitem{diliberti-rogers} I.~Di~Liberti, M.~Rogers, \emph{Topoi with enough points}, arXiv:2403.15338, 2024.
\bibitem{hora-kamio} R.~Hora, Y.~Kamio, \emph{Quotient toposes of discrete dynamical systems}, J. Pure Appl. Algebra 228 (2024), no.~8, 107657; arXiv:2310.02647.
\bibitem{caramello-unification} O.~Caramello, \emph{The unification of Mathematics via Topos Theory}, arXiv:1006.3930, 2010.
\bibitem{caramello-tst} O.~Caramello, \emph{Theories, Sites, Toposes: Relating and studying mathematical theories through topos-theoretic `bridges'}, Oxford University Press, 2017.
\bibitem{giraud} J.~Giraud, \emph{Cohomologie non abélienne}, Grundlehren der math. Wissenschaften~179, Springer, 1971 (Zbl 0226.14011).
\bibitem{mumford} D.~Mumford, Picard groups of moduli problems, in \emph{Arithmetical Algebraic Geometry} (Purdue, 1963), Harper \& Row, New York, 1965, 33--81.
\bibitem{edidin-hu} D.~Edidin, Z.~Hu, The integral Chow rings of the stacks of hyperelliptic Weierstrass points, \emph{Michigan Math.\ J.} 76 (2026), no.~5, 1091--1113; arXiv:2208.00556.
\bibitem{maltsiniotis} G.~Maltsiniotis (ed.), \emph{Alexandre Grothendieck, À la poursuite des champs}, vol.~I, Documents mathématiques~20, Société mathématique de France, 2022.
\bibitem{hua-tangent} M.~Hua, \emph{The tangent line at a marked point on $M_{0,n}/\mathfrak S_{n-1}$, after a table of Grothendieck}, Hexagon preprint 2611.00189, 2026, \url{https://hexagonmath.org/2611.00189}.
\bibitem{bridges-generator} M.~Hua, \emph{Topos-theoretic bridges as a theorem generator}, 2026.
\end{thebibliography}

\end{document}
